\begin{lemma}
\label{lem:ch20-fixed-form-rational-approximation}
Assume that \(\{\theta\}\) is ample.  Every real \(\theta\)-psh model
function is a uniform limit on \(X^{\An}\) of rational \(\theta\)-psh
model functions.  The form \(\theta\) is fixed throughout and need not
itself be semipositive or rational.
\end{lemma}

\begin{proof}
Let \(f=\varphi_D\) be a real \(\theta\)-psh model function and fix
\(\varepsilon>0\).  By
\Cref{prop:ch20-positive-representative}, choose one model \(\mathcal Y\)
determining \(D\) and \(\theta\), and a vertical real Cartier divisor
\(P\) on \(\mathcal Y\), such that \(\theta_{\mathcal Y}+[P]\) is
relatively ample.  Here the cofinal clause is important: the correcting
divisor and the given divisor must live on the same model.  This is also
the positive-representative construction of \cite[Proposition 5.2]{BFJ16}.
For \(0<\delta<1\), set
\[
 D_\delta=(1-\delta)D+\delta P.
\]
Then
\[
 \theta_{\mathcal Y}+[D_\delta]
 =(1-\delta)(\theta_{\mathcal Y}+[D])
       +\delta(\theta_{\mathcal Y}+[P])
\]
is relatively ample, by the sum-of-nef-and-ample assertion in
\Cref{lem:ch20-special-fibre-test}.  Since model functions are continuous
on the compact space \(X^{\An}\), their sup norms are finite.  Choose
\(\delta\) so small that
\(\|\varphi_{D_\delta}-f\|_\infty<\varepsilon/2\).

By the definition of a real Cartier divisor, \(D\) and \(P\) are real
linear combinations of a finite list of integral vertical Cartier divisors
\(F_1,\ldots,F_r\) on this same model.  Write
\(D_\delta=\sum_j a_jF_j\).  The map
\[
 (q_1,\ldots,q_r)\longmapsto
       \theta_{\mathcal Y}+\sum_jq_j[F_j]
\]
is continuous into the finite-dimensional numerical space.  Openness of
its ample cone therefore permits rational numbers \(q_j\) arbitrarily
close to \(a_j\) for which the displayed class remains ample.  Choose
them also so that
\[
 \sum_j|q_j-a_j|\,\|\varphi_{F_j}\|_\infty<\varepsilon/2.
\]
For \(D_\varepsilon=\sum_jq_jF_j\), the rational model function
\(f_\varepsilon=\varphi_{D_\varepsilon}\) is \(\theta\)-psh, and the
triangle inequality gives \(\|f_\varepsilon-f\|_\infty<\varepsilon\).
Only divisor coefficients were rounded; the real class \(\theta\) was
never changed.  Taking \(\varepsilon\) tending to zero proves the result.
\end{proof}

\begin{corollary}
\label{cor:ch20-real-max-uniform-approximation}
If \(\{\theta\}\) is ample and \(f_1,\ldots,f_r\) are real
\(\theta\)-psh model functions, their pointwise maximum is a uniform
limit of rational \(\theta\)-psh model functions.  It need not itself
be a real model function.
\end{corollary}

\begin{proof}
Choose rational approximants to the finitely many \(f_j\) with common
error less than \(\varepsilon\), using
\Cref{lem:ch20-fixed-form-rational-approximation}.  Their maximum is a
rational \(\theta\)-psh model function by \Cref{prop:ch20-model-max}.
The elementary inequality
\[
 \left|\max_j a_j-\max_j b_j\right|\leq\max_j|a_j-b_j|
\]
gives the uniform error bound.  The failure of the last conclusion is
exhibited in \Cref{ex:ch20-real-max-irrational-breakpoint}.
\end{proof}

\begin{example}
\label{ex:ch20-real-max-irrational-breakpoint}
Take \(R=\mathbf C[[\pi]]\), \(K=\mathbf C((\pi))\), and let
\(\rho:\mathcal X\to\mathbf P^1_R\) be the blowup of a rational point
on the special fibre.  Let \(E\) be the exceptional component and \(C\)
the strict transform of that fibre.  They meet transversely in one
residue-rational point, and
\[
 C\cdot E=1,\qquad C^2=E^2=-1,\qquad \mathcal X_s=C+E.
\]
The two self-intersections follow also by intersecting the principal
divisor \(C+E\) with \(C\) and \(E\).  Set
\(\mathcal L=\rho^*\cO_{\mathbf P^1_R}(2)\otimes\cO_{\mathcal X}(-E)\)
and let \(\theta\) be its closed model form.  The two component degrees
are \((1,1)\), so \(\theta\) is semipositive and its generic class has
degree two.  For \(g=\varphi_E\), the curvature form
\(\theta+\ddc g\) is that of \(\rho^*\cO(2)\), with component degrees
\((2,0)\); hence \(g\) is \(\theta\)-psh as well.

The skeleton is the interval \([0,1]\), oriented from \(x_C\) to
\(x_E\), and \(g(s)=s\).  Choose an irrational \(a\in(0,1)\).
Both \(0\) and \(g-a\) are real \(\theta\)-psh model functions, but
\[
                    h=\max\{0,g-a\}
\]
is not a real model function.  Indeed, any real model function is a
finite real linear combination of rational model functions, by
\eqref{eq:ch20-model-functions}.  Each rational model function restricts
to a rational piecewise affine function on this interval
(\Cref{prop:skeleton-model-pa}; see also
\cite[Proposition 3.12 and Corollary 3.13]{BFJ16}).  A finite common
subdivision thus gives only rational possible breakpoints for the real
linear combination.  In contrast, \(h|_{[0,1]}=\max\{0,s-a\}\) has a
genuine breakpoint at the irrational point \(a\).  This contradiction
does not depend on how high a model is chosen.

This also disproves directedness of the real-model minorants of a
continuous obstacle.  Use \(h\) itself as obstacle.  The functions
\(0\) and \(g-a\) are admissible minorants.  Any common upper bound
among those minorants would have to be both at least
\(\max\{0,g-a\}=h\) and at most \(h\), hence equal to the non-model
function \(h\).  There is no such upper bound.  On the other hand,
\Cref{cor:ch20-real-max-uniform-approximation} still applies to \(h\).
Thus neither this example nor the directedness failure contradicts the
maximum stability of genuine psh functions in
\Cref{thm:ch21-stability}.
\end{example}
